\newcommand{\plcode}[1]{%
  \raisebox{-0.15ex}{\tt\textquotesingle}\texttt{{#1}}%
  \raisebox{-0.15ex}{\tt\textquotesingle}%
}
\newcommand{\pldot}{\plcode{.}}
\newcommand{\plcomma}{\plcode{,}}
\newcommand{\ie}{\textit{i.e.}, }

\chapter{Logic Programming}\label{chap:LP}
\chapterauthor{Étienne Payet}

\section{Introduction}

This chapter describes the syntax and semantics for the
\emph{Logic Programming} category of the Termination Competition.
The problems in this category are based on the formalism of
\emph{logic programming with the leftmost selection rule},
extended with some syntactic sugar for lists, tuples, and 
arithmetic operators.
This formalism is also the basic formalism behind the categories
\emph{Logic Programming With Cut}~\cite{Schneider-KampGSST10}
and \emph{Prolog}~\cite{Deransart96,ISO-Prolog95}.
%
For more details on the syntax and semantics of logic programming,
we refer the reader to standard textbooks, for example~\cite{Apt97,Lloyd87}.

\section{Syntax}

Logic programs (LPs) have the following syntax.

\begin{grammar}{LP Syntax}
  \begin{tabular}{lll}
  lp & ::= & query clause+ \\
  query & ::= & {\color{red}\verb+%query:+} IDENTIFIER indicators? {\color{red}\verb+.+} \\
  indicators & ::= & {\color{red}\verb+(+} indicator ( {\color{red}\verb+,+} indicator )* {\color{red}\verb+)+}\\
  indicator & ::= & {\color{red}\verb+i+} $|$ {\color{red}\verb+o+} \\
  clause & ::= & atom {\color{red}\verb+.+} $|$
  atom {\color{red}\verb+:-+} atom ( {\color{red}\verb+,+} atom )* {\color{red}\verb+.+} \\
  atom & ::= & IDENTIFIER $|$ IDENTIFIER tuple \\
  tuple & ::= & {\color{red}\verb+(+} term ( {\color{red}\verb+,+} term )* {\color{red}\verb+)+} \\
  term & ::= & atom $|$ list $|$ tuple $|$ VAR $|$ NUM $|$ term binop term \\
  list & ::= & {\color{red}\verb+[]+} $|$
  {\color{red}\verb+[+} term ( {\color{red}\verb+,+} term )* 
  ({\color{red}\verb+|+} term)? {\color{red}\verb+]+} \\
  binop & ::= & {\color{red}\verb-+-} $|$ {\color{red}\verb+-+} $|$ {\color{red}\verb+*+} $|$ {\color{red}\verb+/+} \\
  \end{tabular}
\end{grammar}

%\bigskip
Here, the red colour is used to highlight the symbols of the syntax,
which are part of the input.
Moreover, IDENTIFIER and VAR denote finite, non-empty strings over
letters, digits, and the underscore character (\verb+_+).
IDENTIFIERs begin with a lowercase letter, whereas VARs begin
with either an uppercase letter or an underscore. Furthermore,
NUM denotes a decimal number.

Additionally, comments are allowed in the input. They can be placed anywhere.
They have two possible forms: they can either start with the symbol \verb+%+
and end with the next newline character, or they can be enclosed between the
symbols \verb+/*+ and \verb+*/+. Every valid LP starts with a comment
of the form \verb+%query: ...+ that specifies the queries of interest, see
Section~\ref{section:LP-competition} below.

The following are examples of valid LPs. They can be found in the
directory \texttt{TPDB/Logic\_Programming} of the TPDB~\cite{tpdb}.

\begin{example}[\texttt{talp\_apt/sum.pl}]\label{ex:lp-sum:syntax}
  % Addition of two natural numbers in Peano arithmetic:
  \phantom{x}
  \begin{verbatim}
%query: sum(o,o,i).

% Addition of two natural numbers in Peano arithmetic.
sum(X,s(Y),s(Z)) :- sum(X,Y,Z).
sum(X,0,X).
  \end{verbatim}
\end{example}

\begin{example}[\texttt{terminweb\_new/append-ooi.pl}]\label{ex:lp-append:syntax}
  This example illustrates the use of the list syntax \verb+[...]+.
  % Append of two lists:
  \begin{verbatim}
%query: app(o,o,i).

% Append of two lists.
app([],X,X).
app([X|Xs],Ys,[X|Zs]) :- app(Xs,Ys,Zs).
  \end{verbatim}
\end{example}

\begin{example}[\texttt{terminweb\_old/balance\_tree.pl}]\label{ex:lp-balance:syntax}
  This example illustrates the use of the binary operator \verb+-+
  and the tuple syntax \verb+(...)+.
  % Balance a binary tree:
  \begin{verbatim}
%query: balance(i,o).

% Balance a binary tree.
balance(T,TB) :-
  balance(T, I-[], [(TB,I-[])|X]-X, Rest-[], Rest-[]).

balance(nil, X-X, A-B, A-B, [(nil,C-C)|T]-T).
balance(tree(L,V,R), IH-IT,
        [(tree(LB,VB,RB),A-D)|H]-[(LB,A-[VB|X]),(RB,X-D)|T],
        HR-TR, NH-NT) :-
  balance(L, IH-[V|IT1], H-T, HR1-TR1, NH-NT1),
  balance(R, IT1-IT, HR1-TR1, HR-TR, NT1-NT).
  \end{verbatim}
\end{example}

\begin{example}[\texttt{talp\_talp/evaluate.pl}]\label{ex:lp-evaluate:syntax}
  This example illustrates the use of the binary operators
  \verb$+$, \verb+-+, and \verb+*+.
  % Arithmetic expression evaluator:
  \begin{verbatim}
%query: myis(o,i).

% Arithmetic expression evaluator.
myis(Z, X) :- evaluate(X, Z).

evaluate(X + Y, Z) :-
        evaluate(X, X1),
        evaluate(Y, Y1),
        add(X1, Y1, Z).
evaluate(X - Y, Z) :-
        evaluate(X, X1),
        evaluate(Y, Y1),
        sub(X1, Y1, Z).
evaluate(X * Y, Z) :-
        evaluate(X, X1),
        evaluate(Y, Y1),
        mult(X1, Y1, Z).
evaluate(X, X) :- myinteger(X).

myinteger(s(X)) :- myinteger(X).
myinteger(0).

add(s(X), Y, s(Z)) :- add(X, Y, Z).
add(0, X, X).

sub(s(X), s(Y), Z) :- sub(X, Y, Z).
sub(X, 0, X).

mult(s(X), Y, R) :- mult(X, Y, Z), add(Y, Z, R).
mult(0, Y, 0).
  \end{verbatim}
\end{example}

\section{Semantics}\label{section:LP-semantics}

We use the same definitions and notation as in Section~\ref{section:TRS-semantics}
for terms and substitutions.
We write $\mathit{Var}(t)$ for the set of variables occurring in a term $t$.
A term $t$ is called \emph{ground} if $\mathit{Var}(t) = \varnothing$.

\subsection{Renaming, Variants, and Unification}
A substitution is called a \emph{renaming} if it is a bijection on the set of variables.
Two terms $s$ and $t$ are said to be \emph{variants} of each other if there exists a
renaming $\rho$ such that $t = s\rho$.

Let $s$ and $t$ be two terms. A substitution $\theta$ is a \emph{unifier} of
$s$ and $t$ if $s\theta = t\theta$. The terms $s$ and $t$ are \emph{unifiable}
if they admit a unifier. A unifier $\theta$ of $s$ and $t$ is a
\emph{most general unifier} (mgu) if every unifier of $s$ and $t$ is an instance
of $\theta$, that is, if for every unifier $\sigma$ there exists a substitution
$\delta$ such that $\sigma = \theta\delta$.
If two terms are unifiable, then they admit a most general unifier,
which is unique up to renaming.

\subsection{Logic Programs}
We assume that the signature $\F$ contains special symbols called
\emph{predicate symbols}.
An \emph{atom} is an expression of the form $p(t_1, \ldots, t_k)$, where
$p$ is a predicate symbol of arity $k \ge 0$ and $t_1, \ldots, t_k$ are
terms. In particular, if $k = 0$, then the atom is simply written $p$.
A \emph{query} is a finite sequence of atoms. We let $\square$ denote
the empty query.
A \emph{clause} is an expression of the form $h \leftarrow b_1,\ldots,b_n$,
where $h,b_1,\ldots,b_n$ are atoms and $n \ge 0$.
The atom $h$ is called the \emph{head} of the clause, and the query
$b_1,\ldots,b_n$ its \emph{body}. The body may be empty, in which case
the clause is called a \emph{fact}. A \emph{logic program} is a finite
set of clauses.

\begin{example}\label{ex:lp-sum:semantics}
  The following logic program has the input syntax of Example~\ref{ex:lp-sum:syntax}:
  \[\left\{
    \mathsf{sum}(x, \mathsf{s}(y), \mathsf{s}(z)) \leftarrow \mathsf{sum}(x, y, z),\
    \mathsf{sum}(x, 0, x) \leftarrow \square
  \right\}\]
\end{example}

\begin{example}\label{ex:lp-append:semantics}
  Consider the logic program 
  \[\left\{
    \mathsf{app}(\texttt{[]}, x, x) \leftarrow \square,\
    \mathsf{app}(\pldot(x,\mathit{xs}), \mathit{ys}, \pldot(x,\mathit{zs}))
    \leftarrow \mathsf{app}(\mathit{xs}, \mathit{ys}, \mathit{zs})
  \right\}\]
  where \texttt{[]} is a constant symbol and \pldot{} is
  a function symbol of arity 2.
  It has the input syntax of Example~\ref{ex:lp-append:syntax}.
  Indeed, the list notation \verb+[...]+ is just syntactic sugar for the 
  Prolog list constructor \pldot. For example,
  the input \verb+[X|Xs]+ is just a shorthand for \pldot\verb+(X,Xs)+
  and the input \verb+[a,b]+ is a shorthand for
  \pldot\verb+(a,+\pldot\verb+(b,[]))+.
\end{example}

\begin{example}\label{ex:lp-balance:semantics}
  Consider the term
  \[\plcomma\left(\mathit{tb},\ -(i,\texttt{[]})\right)\]
  where $\mathit{tb}$ and $i$ are variables,
  \texttt{[]} is a constant symbol, and
  \plcomma{} is a function symbol of arity 2.
  It has the input syntax \verb+(TB,I-[])+ of
  Example~\ref{ex:lp-balance:syntax}.
  Indeed, the tuple notation \verb+(...)+ is just syntactic
  sugar for the Prolog tuple constructor \plcomma, \ie
  the input \verb+(TB,I-[])+ is a shorthand for
  \plcomma\verb+(TB, -(I,[]))+.
\end{example}

\subsection{Left-Derivations, Left-Termination}
Throughout this chapter, we consider SLD-resolution with the
leftmost selection rule. Following the standard terminology
of the termination literature, the resulting derivations are
called \emph{left-derivations}.

Let $Q = a_1, \ldots, a_n$ be a non-empty query and $c$ be a clause.
Let $c' = (h \leftarrow b_1, \ldots, b_m)$ be a variant of $c$,
variable-disjoint from $Q$. If the leftmost atom $a_1$ and the head
$h$ are unifiable with mgu $\theta$, then $Q$ is said to \emph{derive}
the query $Q' = (b_1, \ldots, b_m, a_2, \ldots, a_n)\theta$
in one \emph{left-derivation step}. This step is denoted by
$Q \Rightarrow_c^{\theta} Q'$,
and $c'$ is called its \emph{input clause}.

Let $P$ be a logic program and $Q_0$ be a query. A \emph{left-derivation} of
$P \cup \{Q_0\}$ is a maximal sequence of left-derivation steps
$Q_0 \Rightarrow_{c_1}^{\theta_1} Q_1 \Rightarrow_{c_2}^{\theta_2} \ldots$
where $c_1,c_2, \ldots$ are clauses of $P$ and the \emph{standardization apart}
condition holds, that is, each input clause used is variable-disjoint from the
initial query $Q_0$ and from the mgu's and the input clauses of the previous steps.
%
% A finite left derivation may end up either with the empty query (then,
% it is a \emph{successful} left-derivation) or with a non-empty query
% (then, it is a \emph{failed} left-derivation).
We say that $P \cup \{Q_0\}$ is \emph{left-terminating}, or equivalently
that $Q_0$ is left-terminating w.r.t. $P$, if there is no infinite
left-derivation of $P \cup \{Q_0\}$.

\begin{example}\label{ex:lp-sum:termination}
  The logic program $P$ of Example~\ref{ex:lp-sum:semantics}
  consists of the clauses
  %
  \begin{align*}
    c_1 &: \mathsf{sum}(x, \mathsf{s}(y), \mathsf{s}(z)) \leftarrow \mathsf{sum}(x, y, z) \\
    c_2 &: \mathsf{sum}(x, 0, x) \leftarrow \square
  \end{align*}
  %
  Let $Q_0$ be the query $\mathsf{sum}(x, y, \mathsf{s}(\mathsf{0}))$.
  The only left-derivations of $P \cup \{Q_0\}$ are:
  \begin{itemize}
    \item $Q_0 \Rightarrow_{c_1}^{\theta_1} \mathsf{sum}(x, y_1, \mathsf{0})
    \Rightarrow_{c_2}^{\theta_2} \square$

    where $\theta_1 = \{x_1/x, y/\mathsf{s}(y_1), z_1/\mathsf{0}\}$
    and $\theta_2 = \{x/\mathsf{0}, x_2/\mathsf{0}, y_1/0\}$.
    %
    \item $Q_0 \Rightarrow_{c_2}^{\theta} \square$

    where $\theta = \{x/\mathsf{s}(\mathsf{0}), y/0, x_1/\mathsf{s}(\mathsf{0})\}$.
  \end{itemize}
  (the input clause used at the level $i$ of each derivation is obtained
  by adding the subscript $i$ to all the variables of the corresponding
  program clause).
  Since there is no infinite left-derivation, $P \cup \{Q_0\}$ is left-terminating.
\end{example}


\section{Termination Competition}\label{section:LP-competition}

The Logic Programming category is concerned with the question
``Are all queries of a given set left-terminating?''.

\begin{formalproblem}
  {Logic Programming}
  {A LP $P$ and a set of queries $\mathcal{Q}$}
  {Are all queries in $\mathcal{Q}$ left-terminating w.r.t. $P$?}
  {See \hyperref[scoring_termination]{Termination Scoring}}
\end{formalproblem}

The set of queries $\mathcal{Q}$ is specified in the input
as a comment of the form
\begin{center}
  \texttt{\%query: IDENTIFIER indicators.}
\end{center}
Here, \texttt{IDENTIFIER} is the name of a predicate symbol, and
\texttt{indicators} is an optional sequence %, enclosed in parentheses,
of indicators, each being either \texttt{i} or \texttt{o}, that
specifies the input/output mode of the arguments of the predicate symbol.
The input mode \texttt{i} specifies an argument that must be instantiated
by a ground term, whereas the output mode \texttt{o} imposes no restriction.

\begin{example}\label{ex:lp-sum:competition}
  The logic program $P$ of Example~\ref{ex:lp-sum:semantics}
  has the input syntax of Example~\ref{ex:lp-sum:syntax},
  where the comment
  \begin{center}
    \texttt{\%query: sum(o,o,i).}
  \end{center}
  specifies the set 
  $\mathcal{Q} = \left\{\mathsf{sum}(t_1, t_2, t_3) \mid
  t_1, t_2, t_3 \in \mathcal{T}(\F,\V), \
  t_3 \text{ is ground}\right\}$.
  Every left-derivation step decreases the number of
  occurrences of the symbol $\mathsf{s}$ in the third
  argument of $\mathsf{sum}$, and thus there is no infinite
  left-derivation of $P \cup \{Q_0\}$ for any query $Q_0$
  in $\mathcal{Q}$. Hence, the best output would be \texttt{YES}.
\end{example}

\begin{example}[\texttt{TPDB/Logic\_Programming/Payet\_22/payet-nonloop-1.pl}]\label{ex:lp-payet22:competition}
  \phantom{x}

  \noindent
  Let $P$ be the logic program that consists of the clauses
  \begin{align*}
    c_1 &: \mathsf{p}(\mathsf{0}, y) \leftarrow \mathsf{p}(\mathsf{s}(y), \mathsf{s}(y)) \\
    c_2 &: \mathsf{p}(\mathsf{s}(x), y) \leftarrow \mathsf{p}(x, y)
  \end{align*}
  It has the input syntax:
  \begin{verbatim}
%query: p(i,i).

p(0, Y) :- p(s(Y), s(Y)).
p(s(X), Y) :- p(X, Y).
  \end{verbatim}
  where the comment
  \begin{center}
    \texttt{\%query: p(i,i).}
  \end{center}
  specifies the set of queries
  $\mathcal{Q} = \left\{\mathsf{p}(t_1, t_2) \mid
  t_1 \text{ and } t_2 \text{ are ground terms}\right\}$.
  This set contains the query $Q_0 = \mathsf{p}(\mathsf{0}, \mathsf{0})$,
  and $P \cup \{Q_0\}$ is not left-terminating.
  Indeed, there is an infinite left-derivation of $P \cup \{Q_0\}$ that
  starts as follows
  (for the sake of simplicity, we omit the mgu used at each step):
  \[\mathsf{p}(\mathsf{0}, \mathsf{0}) \Rightarrow_{c_1}
  \mathsf{p}(\mathsf{s}(\mathsf{0}), \mathsf{s}(\mathsf{0})) \Rightarrow_{c_2}
  \mathsf{p}(\mathsf{0}, \mathsf{s}(\mathsf{0})) \Rightarrow_{c_1}
  \mathsf{p}(\mathsf{s}^2(\mathsf{0}), \mathsf{s}^2(\mathsf{0})) \Rightarrow_{c_2}
  \cdots\]
  Hence, the best output would be \texttt{NO}.
\end{example}

If the arity of the predicate symbol \texttt{IDENTIFIER} is $0$
then no indicators are needed and the parentheses are omitted,
as in the following example.

%\newpage
\begin{example}[\texttt{TPDB/Logic\_Programming/talp\_apt/lte.pl}]\label{ex:lp-lte:competition}
  \phantom{x}
  \begin{verbatim}
%query: goal.

even(s(s(X))) :- even(X).
even(0).

lte(s(X),s(Y)) :- lte(X,Y).
lte(0,Y).

goal :- lte(X,s(s(s(s(0))))), even(X).
  \end{verbatim}
  Here, the comment
  \begin{center}
    \texttt{\%query: goal.}
  \end{center}
  specifies the set of queries $\mathcal{Q} = \left\{\mathsf{goal}\right\}$.
\end{example}


\subsection{The Occurs-Check}

The occurs-check is a part of algorithms for unification.
It causes unification of a variable $x$ and a term $t$
to fail if $t$ contains $x$ and $t$ is not equal to $x$.
For efficiency reasons, most
Prolog implementations do not perform the occurs-check, and 
instead they allow the construction of cyclic terms, which are
infinite terms that contain themselves as subterms. This can
lead to different termination behavior, as illustrated by the
following example.
In the \emph{Logic Programming} category of the
Termination Competition, the best output must be computed under
the assumption that the occurs-check is performed.

\begin{example}\label{ex:lp-occurs-check:competition}
  Let $P$ be the logic program that consists of the clauses
  \begin{align*}
    c_1 &: \mathsf{p}(x) \leftarrow \mathsf{q}(x,x) \\
    c_2 &: \mathsf{q}(\mathsf{f}(y),y) \leftarrow \mathsf{p}(y)
  \end{align*}
  It has the input syntax:
  \begin{verbatim}
%query: p(o).
p(X) :- q(X, X).
q(f(Y), Y) :- p(Y).
  \end{verbatim}
  where the comment
  \begin{center}
    \texttt{\%query: p(o).}
  \end{center}
  specifies the set of queries
  $\mathcal{Q} = \left\{\mathsf{p}(t) \mid
  t \in \mathcal{T}(\F,\V)\right\}$.
  \begin{itemize}
    \item Suppose that the occurs-check is performed.
    Let $Q_0$ be any element of $\mathcal{Q}$. Then, $Q_0$ has the form
    $\mathsf{p}(t_0)$ for some term $t_0$.
    We have the left-derivation step
    $Q_0 \Rightarrow_{c_1} \mathsf{q}(t_0,t_0)$
    (we omit the mgu). Let $y'$ be a fresh variable
    and let $c_2'$ be the variant of $c_2$ obtained by
    replacing $y$ with $y'$.
    Resolving the query $\mathsf{q}(t_0,t_0)$ against the clause
    $c_2$ with the input clause $c_2'$ requires unifying $\mathsf{q}(t_0,t_0)$
    with $\mathsf{q}(\mathsf{f}(y'),y')$, which yields the equations
    $t_0=\mathsf{f}(y')$ and $t_0=y'$, hence $y'=\mathsf{f}(y')$.
    As the occurs-check is performed, this unification fails and
    the derivation finitely fails.
    Consequently, as $Q_0$ denotes any element of $\mathcal{Q}$,
    the best output would be \texttt{YES}.

    \item Now, suppose that the occurs-check is not performed.
    Let $Q_0 = \mathsf{p}(\mathsf{f}(z))$. Then, $Q_0$ is an element
    of $\mathcal{Q}$, and we have the left-derivation step
    $Q_0 \Rightarrow_{c_1} \mathsf{q}(\mathsf{f}(z),\mathsf{f}(z))$.
    % (we omit the mgu).
    Resolving the query $\mathsf{q}(\mathsf{f}(z),\mathsf{f}(z))$
    against the clause $c_2$ requires unifying
    $\mathsf{q}(\mathsf{f}(z),\mathsf{f}(z))$ with
    $\mathsf{q}(\mathsf{f}(y),y)$, which yields the equations
    $\mathsf{f}(z)=\mathsf{f}(y)$ and $\mathsf{f}(z)=y$,
    hence $y=\mathsf{f}(y)$.
    %
    As the occurs-check is not performed, the unification succeeds
    by constructing a cyclic term. The remaining query is
    $Q_1 = \mathsf{p}(y)\theta$, where $\theta=\{y/\mathsf{f}(y)\}$,
    \ie $Q_1 = \mathsf{p}(\mathsf{f}(y))$,
    where $y$ satisfies the equation $y=\mathsf{f}(y)$.    
    Resolving $Q_1$ again against $c_1$ and $c_2$ yields a query
    of the same form (up to renaming). Repeating this reasoning
    indefinitely yields an infinite derivation.
    Hence, the best output would be \texttt{NO}.
  \end{itemize}
\end{example}